%%%%%%%%%%%%%%%%%%%%%%%%%%%%%%%%%%%%%%%%%%%%%%%%%%%%%%%%%%%%%
\section{Mathematical Models}\label{sec:math_model}
%%%%%%%%%%%%%%%%%%%%%%%%%%%%%%%%%%%%%%%%%%%%%%%%%%%%%%%%%%%%%
% Rewritten 1 Nov 2026 (overlap split with the companion solver paper):
% same model and notation, condensed, displacement-based solid only.

Every benchmark couples an incompressible Newtonian fluid to an elastic solid.
The fluid occupies $\Omega_F(t)$ and the solid $\Omega_S(t)$, and they meet on the interface $\Gamma_\text{fsi}(t) = \partial\Omega_F(t) \cap \partial\Omega_S(t)$; $\hat{\bb{n}}_F$ and $\hat{\bb{n}}_S$ are the outward unit normals of the two regions, so that $\hat{\bb{n}}_F = -\hat{\bb{n}}_S$ on $\Gamma_\text{fsi}$.
Subscripts $F$ and $S$ mark the region and are dropped where the region is clear, and a subscript $0$ marks the initial configuration.
Following OpenFOAM, the gradient of a vector is $(\bb{\nabla}\bb{a})_{ij} = \partial a_j/\partial x_i$, the transpose of the usual continuum-mechanics convention, and $\bb{\nabla}_0$ differentiates with respect to the initial coordinates.
Body forces are per unit mass.

%%--------------------------------------------------------------------------------------------------------------------%%
\subsection{Solid model} \label{sec:solid_model}
%%--------------------------------------------------------------------------------------------------------------------%%

The solid is described in total Lagrangian form with the displacement $\bb{D}_S$ from the initial configuration $\Omega_{S,0}$ as the unknown.
With the deformation gradient $\bb{F} = \bb{I} + (\bb{\nabla}_0 \bb{D}_S)^T$, $J = \det\bb{F}$, and the first Piola--Kirchhoff stress $\bb{P}_S = J\bb{\sigma}_S\bb{F}^{-T}$, where $\bb{\sigma}_S$ is the Cauchy stress, the balance of linear momentum over any part of the initial configuration reads
\begin{equation}
    \int_{\Omega_{S,0}} \rho_{S,0} \frac{\partial^2 \bb{D}_S}{\partial t^2} \, \text{d}\Omega_{S,0}
    =
    \oint_{\Gamma_{S,0}} \hat{\bb{n}}_{S,0} \cdot \bb{P}_S^T \, \text{d}\Gamma_{S,0}
    +
    \int_{\Omega_{S,0}} \rho_{S,0} \bb{b}_S \, \text{d}\Omega_{S,0},
    \label{eqn:solid_momentum_TL}
\end{equation}
with the time derivative taken at fixed material coordinates, so that $\bb{v}_S = \partial\bb{D}_S/\partial t$.
The stress comes from linear elasticity or from a neo-Hookean or Saint Venant--Kirchhoff hyperelastic law.
Where the strains and rotations are small, the geometrically linear form is used instead: $\bb{F} \approx \bb{I}$ and $\bb{P}_S \approx \bb{\sigma}_S$ in Equation~\eqref{eqn:solid_momentum_TL}, with Hooke's law and Lam\'e parameters $\mu_S$ and $\lambda_S$.
The solid boundaries carry prescribed displacements, prescribed tractions or symmetry conditions, besides the interface conditions of Section~\ref{sec:fsi_interface}.

%%--------------------------------------------------------------------------------------------------------------------%%
\subsection{Fluid model}
\label{sec:fluid_governing_eqn}
%%--------------------------------------------------------------------------------------------------------------------%%

The fluid is laminar, isothermal, incompressible and Newtonian, and is described on the moving domain $\Omega_F(t)$ in arbitrary Lagrangian--Eulerian (ALE) form, with $\bb{v}^\Gamma$ the velocity of the domain.
The domain follows the solid on $\Gamma_\text{fsi}$; inside $\Omega_F$ its motion is a numerical choice (Section~\ref{sec:fluid_mesh_motion}) that must satisfy the geometric conservation law
\begin{equation}
    \frac{d}{dt} \int_{\Omega_F} \text{d}\Omega_F - \oint_{\Gamma_F} \hat{\bb{n}}_F \cdot \bb{v}^\Gamma \, \text{d}\Gamma_F = 0 .
    \label{eqn:gcl}
\end{equation}
With Equation~\eqref{eqn:gcl}, mass conservation over a moving control volume reduces to $\oint_{\Gamma_F} \hat{\bb{n}}_F \cdot \bb{v}_F \, \text{d}\Gamma_F = 0$, and momentum conservation is
\begin{equation}
    \frac{d}{dt} \int_{\Omega_F} \bb{v}_F \, \text{d}\Omega_F
    + \oint_{\Gamma_F} \hat{\bb{n}}_F \cdot (\bb{v}_F - \bb{v}^\Gamma)\, \bb{v}_F \, \text{d}\Gamma_F
    = \oint_{\Gamma_F} \hat{\bb{n}}_F \cdot \bb{\sigma}_F^* \, \text{d}\Gamma_F
    + \int_{\Omega_F} \bb{b}_F \, \text{d}\Omega_F,
    \label{eqn:fluid_momentum_ale}
\end{equation}
with the kinematic stress $\bb{\sigma}_F^* = \nu\left[\bb{\nabla}\bb{v}_F + (\bb{\nabla}\bb{v}_F)^T\right] - p_F \bb{I}$, where $p_F$ is the kinematic pressure (pressure over density), $\nu$ the kinematic viscosity, and $\bb{\sigma}_F = \rho_F \bb{\sigma}_F^*$ the Cauchy stress.
The fluid boundaries are inlets, outlets, no-slip walls and symmetry planes; the conditions of each benchmark are those of its solids4foam tutorial.

%%--------------------------------------------------------------------------------------------------------------------%%
\subsection{Interface conditions}
\label{sec:fsi_interface}
%%--------------------------------------------------------------------------------------------------------------------%%

On $\Gamma_\text{fsi}$ the fluid sticks to the solid, which also moves the fluid domain, and the tractions balance:
\begin{equation}
    \bb{v}_F = \bb{v}^\Gamma = \frac{\partial \bb{D}_S}{\partial t},
    \qquad
    \rho_F \, \hat{\bb{n}}_F \cdot \bb{\sigma}_F^* + \hat{\bb{n}}_S \cdot \bb{\sigma}_S = \bb{0}
    \qquad \text{on } \Gamma_\text{fsi}(t).
    \label{eqn:kinematic_interface}
\end{equation}
With consistent initial conditions, the kinematic condition keeps the interface free of gaps and overlaps.
